\documentclass[11pt,a4paper]{article}
% Shared article style for handouts.
\usepackage[fontset=none]{ctex}
\usepackage{fontspec}
\setmainfont{DejaVu Serif}
\setsansfont{DejaVu Sans}
\setmonofont{DejaVu Sans Mono}
\setCJKmainfont{Noto Serif CJK SC}[BoldFont={Noto Serif CJK SC Bold}]
\setCJKsansfont{Noto Sans CJK SC}[BoldFont={Noto Sans CJK SC Bold}]
\setCJKmonofont{Noto Sans Mono CJK SC}

\usepackage[margin=1in]{geometry}
\usepackage{graphicx}
\usepackage{enumitem}
\usepackage{amsmath,amssymb}
\usepackage{booktabs,array}
\usepackage{xcolor}
\usepackage{listings}
\usepackage{tcolorbox}
\tcbuselibrary{breakable}
\usepackage{hyperref}

\definecolor{hcPrimary}{RGB}{20,72,110}
\definecolor{hcTipBG}{RGB}{239,246,252}
\definecolor{hcTipFrame}{RGB}{80,126,160}
\definecolor{hcWarnBG}{RGB}{253,244,238}
\definecolor{hcWarnFrame}{RGB}{183,84,52}
\definecolor{hcCodeBG}{RGB}{247,248,250}
\definecolor{hcCodeFrame}{RGB}{208,214,220}

\hypersetup{colorlinks=true,linkcolor=hcPrimary,urlcolor=hcPrimary,citecolor=hcPrimary}
\setlist{nosep}

\newtcolorbox{tipbox}{
  colback=hcTipBG,
  colframe=hcTipFrame,
  title=提示,
  fonttitle=\bfseries,
  boxrule=0.6pt,
  breakable
}

\newtcolorbox{warnbox}{
  colback=hcWarnBG,
  colframe=hcWarnFrame,
  title=警告,
  fonttitle=\bfseries,
  boxrule=0.6pt,
  breakable
}

\lstset{
  basicstyle=\ttfamily\footnotesize,
  backgroundcolor=\color{hcCodeBG},
  frame=single,
  rulecolor=\color{hcCodeFrame},
  numbers=left,
  numberstyle=\tiny,
  breaklines=true,
  showstringspaces=false,
  tabsize=2
}


\title{<Lecture Title>}
\author{<Course Name>}
\date{\today}

\graphicspath{{figures/}}

\begin{document}
\maketitle

\section{本讲导读}

% 先交代：为什么有这一讲？它和前后章节是什么关系？
% 不要一上来堆术语，先给读者位置感。
<Why this lecture exists. What problem does it solve in the course map?>

\section{Learning Objectives}

% 建议 4--8 条。写“读完这一讲，你会……”，不要写成训诫式清单。
\begin{itemize}
  \item <Objective 1>
  \item <Objective 2>
  \item <Objective 3>
\end{itemize}

\section{为什么这一讲值得学}

% 这一节回答主问题。先立主线，再给定义。
% 如果这是第一次引入一个概念，先讲轮廓和作用；细节可以后面螺旋式回来加深。
<Main motivation and core question.>

\section{核心概念与对象}

% 推荐组织方式：对象、关系、原因、边界。
% 尽量让读者感觉是在接近“对象本身的结构”，而不是在接受作者立场。
<Core concept explanation here.>

\section{关键图示或案例}

% 至少放一个关键图、一张表或一个现实系统例子。
% 图示的作用是帮助建立结构感，不是装饰。
<Key figure / table / system example here.>

\begin{tipbox}
% 用于建立直觉、给观察方法、降低阅读焦虑。
% 也可以提示：“这一遍先抓轮廓，后面会回来加深。”
<Use this box for study hints, intuition, or a spiral-learning callback.>
\end{tipbox}

\begin{warnbox}
% 用于指出常见误解、抽象失效条件、边界条件或容易混淆的概念。
<Use this box for risks, misconceptions, or boundary conditions.>
\end{warnbox}

\section{本讲小结}

% 回收本讲主线：读者现在应该能回答什么问题？
% 尽量把本讲内容重新接回课程地图。
<Short recap and bridge to later chapters.>

\section{练习题}

% 练习题优先要求：解释、画图、分类、对比、分析。
% 尽量让题目在训练一种观察方式，而不只是检查记忆。
\begin{enumerate}
  \item <Question 1>
  \item <Question 2>
  \item <Question 3>
\end{enumerate}

\end{document}
